% !TEX root = ../../main.tex

\section{Spring Boot 服务端核心模块实现}

Spring Boot 服务端是系统业务中心，主要负责用户认证、资源管理、文件管理、句子视频识别编排、语义修正调用、模型版本管理和后台管理端接口支持。
与移动端和管理端相比，服务端的重点在于统一数据口径、完成权限校验、封装跨服务调用并维护业务状态。

服务端核心业务接口如表~\ref{tab:chapter05-backend-interfaces} 所示。

\WhuInterfaceTable{服务端核心业务接口}{tab:chapter05-backend-interfaces}{
认证管理 & POST & /app/auth/login & 普通用户登录认证 \\
管理员认证 & POST & /admin/auth/login & 后台管理员登录认证 \\
分类管理 & GET/POST & /admin/sign/categories & 手势分类查询与维护 \\
资源管理 & GET/POST & /admin/sign/resources & 手势资源查询与维护 \\
样本管理 & GET/POST & /admin/sign/samples & 手势样本查询与维护 \\
视频识别 & POST & /app/video-recognition & 上传句子视频并组织识别结果 \\
语义修正 & POST & /app/semantic-correction & 根据识别词序列生成语义修正结果 \\
文件上传 & POST & /file/upload & 上传封面、视频、样本或模型文件 \\
模型版本 & GET/POST & /admin/model-versions & 模型版本登记与查询 \\
模型发布 & POST & /admin/model-versions/publish & 发布模型并联动识别服务加载 \\
}

\subsection{服务端分层结构实现}

服务端采用控制器层、服务层、持久层和外部服务客户端相结合的结构。
控制器层负责接收移动端和管理端 HTTP 请求，服务层负责业务规则处理和跨服务编排，持久层基于 MyBatis 完成数据库访问，外部服务客户端负责调用 Python 手势识别服务和 DeepSeek 服务。

当前服务端控制器按业务模块划分，主要包括 AppAuthController、AppUserController、AppRecentPracticeController、SignCategoryController、SignResourceController、SignSampleController、SignModelVersionController、SignVideoRecognitionController、SignSemanticCorrectionController、FileUploadController 等。
该划分方式与系统功能模块保持一致，使移动端、后台管理端和识别相关接口具有较清晰的入口。

服务端接口返回结构尽量保持统一，前端在调用时只需要关注业务数据和错误提示。
对于参数错误、登录状态失效、数据库异常、对象存储异常和识别服务异常，服务端通过统一异常处理机制转换为前端可理解的响应信息，避免直接暴露底层异常。
通过这种分层方式，服务端能够在保持接口结构清晰的同时，降低业务逻辑、数据库访问和外部服务调用之间的耦合。

\subsection{用户认证与登录态维护实现}

服务端同时支持普通用户和管理员两类认证场景。
普通用户认证面向 HarmonyOS 移动端，管理员认证面向 Vue 后台管理端。
两类登录流程在业务角色上相互区分，但基本思路一致：登录成功后生成 token，并将 token 与用户身份信息写入 Redis；
后续请求携带 token 时，服务端根据 Redis 中的登录态判断请求是否有效。

采用 Redis 保存登录态的原因在于，登录状态属于高频读取、短时效和可过期数据，不适合每次请求都访问 MySQL 进行校验。
服务端在用户退出登录或 token 过期后移除对应登录态，使客户端必须重新登录。
通过这种方式，系统既保证了接口访问控制，也降低了数据库认证压力。

在移动端场景中，用户登录成功后，移动端保存 token 和必要的用户信息，并在后续请求中携带 token。
服务端根据 token 查询 Redis 中的用户登录态，并将用户身份传递给后续业务逻辑。
对于后台管理端，管理员登录后同样获得 token，但服务端在处理后台接口时会区分管理员身份和权限场景，避免普通用户访问管理接口。

退出登录时，客户端主动调用退出接口，服务端删除 Redis 中对应登录态。
若 token 已经过期或不存在，服务端会向客户端返回登录失效信息，前端再跳转到登录页面或要求用户重新登录。
该流程使普通用户端和后台管理端都能够形成较完整的认证闭环。

\subsection{手语分类与手语资源服务实现}

手语分类和手语资源服务为移动端训练内容展示提供数据支撑，也为后台管理端资源维护提供接口。
分类服务主要负责分类列表查询、新增、修改、删除和状态控制；
资源服务主要负责手语资源的分页查询、详情查询、新增、修改、删除、启用禁用和资源文件关联等操作。

对于移动端，服务端主要提供已启用分类和已启用资源的查询接口。
移动端通过这些接口构建首页分类入口、资源列表和训练详情页面。
由于移动端面向普通用户，接口返回内容需要避免暴露过多后台维护字段，而是重点返回资源名称、分类信息、封面地址、示范资源地址和可训练状态等展示所需信息。

对于后台管理端，服务端提供更完整的分类和资源维护接口。
管理员可以在后台新增或修改分类，维护资源编码、资源名称、所属分类、封面资源、动作示范资源和启用状态等信息。
服务端在处理这些操作时需要完成参数校验、资源唯一性检查和数据库更新，保证移动端展示的数据具有一致性。

手语资源服务还承担连接训练流程和识别服务的作用。
资源编码或资源标签需要与 Python 识别服务中的标签体系保持对应关系，使移动端选择的训练资源能够与识别模型输出结果进行匹配。
因此，服务端在资源维护过程中不仅保存展示字段，也需要维护资源编码等与识别链路相关的关键字段。

\subsection{句子视频识别服务实现}

句子视频识别服务用于支持移动端上传手语句子视频，并获得词级识别结果、中文映射结果和后续语义修正结果。
该功能不由移动端直接调用 Python 识别服务完成，而是通过 Spring Boot 服务端进行统一编排。
这样可以避免移动端直接感知多个后端服务的调用细节，也便于服务端统一处理文件保存、接口调用、结果组织和异常提示。

用户在移动端选择本地手语视频后，移动端将视频文件上传到 Spring Boot 服务端。
服务端接收视频后，可以先将视频保存到对象存储或临时目录，再调用 Python 手势识别服务进行视频识别。
Python 服务读取视频帧，对视频中的动作片段进行关键点提取、特征构造和模型预测，生成词级识别结果。

与实时识别不同，句子视频识别需要从完整视频中提取动作序列，并尽可能生成连续词语结果。
Python 服务可以根据视频帧序列逐段识别手势动作，并输出原始标签序列或候选词序列。
由于模型识别结果可能存在重复、遗漏或顺序不够自然的问题，服务端还需要对识别结果进行初步整理。

Spring Boot 服务端在收到 Python 服务返回的结果后，负责将标签结果映射为中文词语，并组织原始识别结果、中文映射结果和待修正文本。
这样可以让系统保留模型原始输出，同时也为后续语义修正提供输入。
移动端最终展示时，可以将这些结果分层呈现给用户。

\subsection{DeepSeek 语义修正服务实现}

DeepSeek 语义修正服务用于将词级识别结果整理为更自然的中文表达。
由于当前系统定位为有限词表和受控场景下的手语辅助识别，模型输出往往是若干词语或短语序列，而不是完全符合汉语语序的自然句。
服务端通过调用 DeepSeek 语义修正服务，对识别结果进行语义组织和表达优化。

服务端在调用语义修正服务时，需要构造清晰的输入文本。
该输入通常包括识别出的中文词序列、任务说明和输出约束。
为了避免语义修正过度发挥，服务端应尽量要求模型基于已有识别结果进行整理，不随意添加识别结果中不存在的核心信息。
这样可以在提升可读性的同时，降低语义修正引入错误内容的风险。

语义修正服务返回结果后，Spring Boot 服务端将其与原始识别结果一起封装返回移动端。
移动端展示时，可以同时显示原始词序列和修正后的句子，使用户能够理解系统识别过程，也可以对比修正前后的表达差异。
该设计有助于提高结果透明度，避免用户只看到一个不可解释的最终句子。

需要注意的是，DeepSeek 在系统中只承担语义后处理职责，不直接参与视频视觉识别，也不替代手势识别模型。
视觉识别结果仍由 Python 手势识别服务生成，DeepSeek 只在识别候选结果的基础上完成语言表达整理。
通过这种分层处理方式，系统能够在保留识别链路可解释性的同时，提高句子结果的可读性。

\subsection{文件上传与资源访问实现}

系统中的头像、封面、动作示范资源、样本文件、视频文件和模型文件等都属于非结构化资源，不适合直接写入数据库。
服务端通过文件上传接口对接 MinIO 对象存储，将文件保存到对象存储中，并在数据库中记录文件访问地址和业务关联关系。
这样既可以避免数据库存储大文件，也便于后续扩展资源访问和文件管理能力。

文件上传模块通常由 FileUploadController 接收前端上传请求，再由文件服务完成文件名生成、存储路径组织、MinIO 上传和访问地址返回。
对于普通图片和视频资源，服务端返回的地址可以直接被移动端或后台管理端使用；
对于模型文件和样本文件，则需要与样本管理或模型版本管理模块建立业务关联。

在手势资源管理场景中，管理员可以上传资源封面、动作示范视频或数字人资源地址，服务端保存文件访问路径后，移动端即可通过资源接口获取并展示对应内容。
在句子视频识别场景中，移动端上传的视频文件可以由服务端临时保存或转发给 Python 识别服务处理。
对于训练样本和模型文件，服务端需要维护其业务归属和版本关系，避免资源文件与数据库记录脱节。

通过文件上传和资源访问模块，系统将结构化业务数据与非结构化文件资源分离存储。
MySQL 保存资源元数据、业务状态和关联关系，MinIO 保存具体文件内容，服务端则负责在二者之间建立映射。
这种设计有利于提升系统资源管理能力，也为后续更换对象存储或扩展文件类型提供了空间。

\subsection{关键接口调用说明}

HearBridge 系统的核心功能依赖多端之间的接口调用。
移动端主要调用普通用户认证、手势分类、手势资源、最近练习和视频识别等接口；
后台管理端主要调用管理员认证、分类管理、资源管理、样本管理、训练管理和模型版本管理等接口；
Spring Boot 服务端在部分业务中还需要调用 Python 手势识别服务和 DeepSeek 语义修正服务。

在实时手势训练场景中，移动端主要通过 WebSocket 与 Python 手势识别服务通信，持续发送图像帧并接收识别结果。
该链路强调连续帧传输和即时反馈，因此不适合通过普通 HTTP 接口反复请求完成。
在句子视频识别场景中，移动端通过 HTTP 上传视频文件，由 Spring Boot 服务端统一编排文件处理、Python 服务调用和语义修正调用。

在后台模型维护场景中，管理员通过后台管理端发起样本查询、模型训练和模型发布等请求。
Spring Boot 服务端负责将这些请求转化为业务操作，并在需要时调用 Python 识别服务完成特征转换、模型训练或模型重载。
模型发布成功后，移动端实时识别和视频识别功能才能使用当前发布模型。

接口调用过程中需要保持统一的数据格式和错误语义。
普通业务接口应返回统一响应结构，识别服务调用失败、文件上传失败、模型加载失败等异常需要被转换为前端可理解的提示信息。
通过统一接口封装和异常处理，系统能够降低多端联调复杂度，提高整体功能稳定性。
